\section{Trace-aware residual certification and exact completion}

The spectral bridge in the preceding section used an operator-norm residual as
if every unresolved coefficient could move independently by the same amount.
That is rigorous but unnecessarily pessimistic once a nontrivial range basis
has already been captured.  The residual has considerably more structure.

Let \(U\in\mathbb R^{n\times r}\) have orthonormal columns and define
\[
B=U^\top A,
\qquad
F=(I-UU^\top)A.
\]
Then
\[
A=UB+F,
\qquad
U^\top F=0.
\]

\begin{theorem}[Exact core--residual decomposition]
\[
\boxed{
A^\top A
=
B^\top B+F^\top F.
}
\]
In particular the unresolved spectral term
\[
R=F^\top F
\]
is positive semidefinite.
\end{theorem}

\begin{proof}
Expand \((UB+F)^\top(UB+F)\).  The two cross terms vanish because
\(U^\top F=0\), while \(U^\top U=I\).
\end{proof}

This changes the certification problem.  Instead of an arbitrary Hermitian
error, the missing term is a positive residual with measurable spectral
moments.  Define
\[
t=\operatorname{tr}R,
\qquad
v=\operatorname{tr}(R^2),
\qquad
\rho=\|R\|_2.
\]
Then
\[
\boxed{
 t=\|F\|_F^2,
 \qquad
 v=\|FF^\top\|_F^2,
 \qquad
 \rho=\|F\|_2^2.
}
\]
The first quantity is the exact missing mean of the Gaussian quadratic form;
the second records how concentrated the residual spectrum is; the third is its
largest possible coefficient.

\subsection{A trace--Frobenius--cap MGF envelope}

Set
\[
C=B^\top B,
\qquad
H=C+R,
\qquad
c=\lambda_{\max}(C).
\]
For \(0<2\theta(c+\rho)<1\), put
\[
\alpha=(1-2\theta c)^{-1}.
\]

\begin{theorem}[Trace--Frobenius--cap log-MGF certificate]
For the Gaussian quadratic form with coefficient matrix \(H\),
\[
K_H(\theta)
=-\frac12\log\det(I-2\theta H)
\]
satisfies
\[
\boxed{
K_H(\theta)
\le
K_C(\theta)
+
\theta\alpha t
+
\frac{\theta^2\alpha^2v}
     {1-2\theta\alpha\rho}.
}
\]
\end{theorem}

\begin{proof}
Factor the determinant through the known core:
\[
I-2\theta(C+R)
=
(I-2\theta C)^{1/2}
(I-2\theta T)
(I-2\theta C)^{1/2},
\]
where
\[
T=(I-2\theta C)^{-1/2}
R
(I-2\theta C)^{-1/2}.
\]
Since
\[
\|(I-2\theta C)^{-1/2}\|_2^2=\alpha,
\]
we have
\[
\operatorname{tr}T\le\alpha t,
\qquad
\operatorname{tr}(T^2)\le\alpha^2v,
\qquad
\|T\|_2\le\alpha\rho.
\]
For every eigenvalue \(x\) of \(T\),
\[
-\frac12\log(1-2\theta x)
\le
\theta x
+
\frac{\theta^2x^2}
     {1-2\theta\alpha\rho}.
\]
The scalar inequality follows from the power series of \(-\log(1-u)\) and a
geometric bound on all terms of degree at least two.  Summing over the
eigenvalues of \(T\) proves the result.
\end{proof}

\begin{corollary}[Trace-aware Chernoff upper tail]
For \(Q=Z^\top HZ\) with \(Z\sim N(0,I)\),
\[
\boxed{
\Pr(Q\ge q)
\le
\inf_{0<2\theta(c+\rho)<1}
\exp\{\widehat K_H(\theta)-\theta q\},
}
\]
where \(\widehat K_H\) is the preceding MGF upper bound.
\end{corollary}

A complementary lower endpoint is immediate from
\[
H=C+R\succeq C.
\]
Using the same normal vector,
\[
Z^\top HZ\ge Z^\top CZ
\]
almost surely, and therefore
\[
\boxed{
\Pr(Z^\top CZ\ge q)
\le
\Pr(Z^\top HZ\ge q).
}
\]
The lower endpoint can be evaluated by K-COS on the captured core spectrum,
while the upper endpoint is obtained from the optimized trace-aware Chernoff
bound.

\subsection{Finite stress test}

The previous coefficient-wise Weyl interval was intentionally conservative.
On the same full-rank, 50-dimensional stress test used in the preceding
section, its \(10^{-8}\) tail interval remained thousands of decades wide
until essentially the complete range had been captured.  Preserving the exact
residual trace and higher residual moments changes that behavior substantially.

\begin{table}[H]
\centering
\caption{Tail-bracket width at a true probability of \(10^{-8}\).  The
trace/moment certificate is still conservative, but it is orders of magnitude
narrower than coefficient-wise Weyl inflation before full range capture.}
\begin{tabular}{rrrr}
\toprule
Krylov depth & basis dimension & v0.7 Weyl width & v0.8 moment width\\
\midrule
4 & 24 & 90,097 & 11.86\\
5 & 30 & 16,511 & 9.76\\
6 & 36 & 3,747 & 7.67\\
7 & 42 & 2,275 & 5.70\\
8 & 48 & 1,165 & 3.39\\
\bottomrule
\end{tabular}
\end{table}

\begin{figure}[H]
\centering
\includegraphics[width=.84\linewidth]{figure_16_trace_moment_vs_weyl.png}
\caption{Incomplete-range tail certification before exact completion.  The
trace/moment bound preserves the amount and concentration of unresolved
spectral mass instead of shifting every coefficient by one worst-case radius.}
\label{fig:trace-moment-vs-weyl}
\end{figure}

A separate random implementation audit used 120 random factor/range pairs and
four admissible positive tilts per pair, for 480 MGF comparisons.  Both the
trace/cap and trace--Frobenius--cap envelopes dominated the exact log MGF in
every tested case.  The trace--Frobenius--cap form was at least as tight as the
trace/cap form in approximately 83.5\% of those comparisons.  This finite
sweep audits the implementation; the theorem above supplies the universal
argument.

\subsection{Residual rank from moments}

The same invariants determine when it is better to stop approximating the
missing spectrum and complete it exactly.  Let
\[
s=\operatorname{rank}F=\operatorname{rank}R.
\]
If the nonzero eigenvalues of \(R\) are \(r_1,\ldots,r_s\), then
\[
t=\sum_i r_i,
\qquad
v=\sum_i r_i^2,
\qquad
r_i\le\rho.
\]

\begin{theorem}[Residual-rank bracket]
\[
\boxed{
\left\lceil
\max\left\{
\frac{t}{\rho},
\frac{t^2}{v}
\right\}
\right\rceil
\le
\operatorname{rank}F
\le
n-r.
}
\]
The zero-residual case is interpreted separately as rank zero.
\end{theorem}

\begin{proof}
The inequality \(t\le s\rho\) gives \(s\ge t/\rho\).  Cauchy--Schwarz gives
\[
t^2
=
\left(\sum_{i=1}^s r_i\right)^2
\le
s\sum_{i=1}^s r_i^2
=sv,
\]
so \(s\ge t^2/v\).  Finally, every column of \(F\) lies in
\(\operatorname{range}(U)^\perp\), whose dimension is \(n-r\).
\end{proof}

Whenever the integer lower bound equals \(n-r\), the missing rank is certified
exactly without computing a residual SVD.

\begin{figure}[H]
\centering
\includegraphics[width=.84\linewidth]{figure_17_residual_rank_completion.png}
\caption{Residual-rank certification in the full-rank stress test.  From depth
3 onward the moment lower bound meets the orthogonal-complement upper bound.
At depth 8 exactly two residual directions remain.}
\label{fig:residual-rank-completion}
\end{figure}

\subsection{Exact residual completion}

Suppose a compact factorization of the residual is
\[
F=WD,
\]
where \(W\) has orthonormal columns.  Since \(F\) lies in the orthogonal
complement of \(U\),
\[
U^\top W=0.
\]

\begin{theorem}[Exact residual completion]
\[
\boxed{
A
=
[U\;W]
\begin{bmatrix}
B\\D
\end{bmatrix}.
}
\]
Consequently, the nonzero singular values of \(A\) are exactly the singular
values of the stacked completion core
\[
\mathcal C_{\rm complete}
=
\begin{bmatrix}
B\\D
\end{bmatrix}.
\]
\end{theorem}

This identity admits a column-streamed implementation.  For each residual
column block, project against the residual basis accumulated so far, append
only genuinely new directions through rank-revealing QR or SVD, and extend the
coordinate matrices of earlier blocks with zero rows.  New residual directions
are orthogonal to every earlier residual block by construction, so no earlier
data block must be revisited.

In the deterministic stress test, exact residual completion reproduced the
full singular spectrum at relative errors between approximately
\(2\times10^{-16}\) and \(6\times10^{-16}\) for all tested incomplete bases.
The resulting \(10^{-8}\) tail differed from the full-spectrum reference by
less than \(10^{-14}\) decades.  At depth 8 the candidate range had dimension
48, the missing rank was certified as exactly two, and appending those two
residual directions closed the spectrum without another approximate Krylov
stage.

\subsection{Hybrid stopping rule}

The resulting production logic is not ``approximate forever.''  It is a
certified handoff between two regimes:
\[
\boxed{
\begin{array}{c}
\text{streamed / pruned Krylov}\\[2mm]
\downarrow\\[-1mm]
\text{residual moments }(t,v,\rho)\\[2mm]
\downarrow\\[-1mm]
\begin{cases}
\text{continue Krylov}, & \text{residual range still large},\\
\text{exact streamed residual completion}, & \text{completion rank affordable}.
\end{cases}
\end{array}
}
\]

The trace-aware tail certificate remains useful as a rigorous intermediate
bound and as a stopping diagnostic.  Once only a small number of residual
directions remain, exact completion is mathematically preferable to forcing a
worst-case tail envelope to become artificially sharp.
